% GENERATED FILE. Edit canonical lessons, then run npm run generate:books.
% canonical-source-hash: fnv1a64:91bbd094bcb5a226
% canonical-lessons: GE-C522-ausleihen, GE-C522-blaettern, GE-C522-lesezeichen, GE-C522-verlag, GE-C522-uebersetzung

\chapter{Bücher leihen}
\label{ch:ge-bucher-leihen}
\hlchaptermodality{522}

\begin{chapteropening}
\textbf{By the end of this chapter:} \emph{I can talk about borrowing and publishing books.}

Complete the last lesson of chapter 522: I can talk about borrowing and publishing books.
\end{chapteropening}

\section[ausleihen]{ausleihen — to borrow; to lend}

\label{lesson:GE-C522-ausleihen}



\begin{quote}
Before the new one: say the German for to rhyme, then the German for writer.
\end{quote}

\subsection*{You'll want to know: ausleihen}
\textbf{ausleihen} — \textquotedblleft{}to borrow; to lend\textquotedblright{}.

\begin{quote}
In der Bibliothek habe ich drei Bücher ausgeliehen.
\end{quote}

\begin{quote}
\emph{I borrowed three books from the library.}
\end{quote}

Separable: \emph{lieh aus, hat ausgeliehen}. \emph{Sich etwas ausleihen} — to borrow; \emph{jemandem etwas ausleihen} — to lend.

\begin{cousinweb}
\emph{ausleihen} is built on \emph{aus} \textquotedblleft{}out\textquotedblright{} and \emph{leihen} \textquotedblleft{}to lend\textquotedblright{}, inherited from Old High German \emph{līhan}, from Proto-Germanic \emph{līhwaną} (EtymDB).
\end{cousinweb}

\begin{grammarlens}[title={What you've built}]
One more word for telling, explaining and describing at B1.
\end{grammarlens}

\subsection*{Your turn}
\begin{itemize}
\raggedright
  \item \emph{Say it:} \emph{ausleihen}
  \item \emph{Say it:} \emph{ausleihen}, once more
  \item \emph{Say it:} a sentence of your own with \emph{ausleihen}
  \item \emph{Recall:} say the German for to rhyme, then the German for writer, then say \emph{ausleihen} again
\end{itemize}

\begin{tcolorbox}[breakable,colback=teal!4,colframe=teal!35!black,title={Before you move on}]
What does ausleihen mean? (\textquotedblleft{}To borrow; to lend\textquotedblright{}.) Say it once more.
\end{tcolorbox}
\section[blättern]{blättern — to leaf (through)}

\label{lesson:GE-C522-blaettern}



\begin{quote}
Before the new one: say the German for writer, then the German for to borrow.
\end{quote}

\subsection*{You'll want to know: blättern}
\textbf{blättern} — \textquotedblleft{}to leaf (through)\textquotedblright{}.

\begin{quote}
Im Wartezimmer blättert sie in einer Zeitschrift.
\end{quote}

\begin{quote}
\emph{In the waiting room she leafs through a magazine.}
\end{quote}

\emph{In einem Buch blättern} — to leaf through a book.

\begin{cousinweb}
\emph{blättern} is built on \emph{Blatt} \textquotedblleft{}leaf, page\textquotedblright{}, inherited from Old High German \emph{blat}, from Proto-Germanic \emph{bladą} (EtymDB).
\end{cousinweb}

\begin{grammarlens}[title={What you've built}]
One more word for telling, explaining and describing at B1.
\end{grammarlens}

\subsection*{Your turn}
\begin{itemize}
\raggedright
  \item \emph{Say it:} \emph{blättern}
  \item \emph{Say it:} \emph{blättern}, once more
  \item \emph{Say it:} a sentence of your own with \emph{blättern}
  \item \emph{Recall:} say the German for writer, then the German for to borrow, then say \emph{blättern} again
\end{itemize}

\begin{tcolorbox}[breakable,colback=teal!4,colframe=teal!35!black,title={Before you move on}]
What does blättern mean? (\textquotedblleft{}To leaf (through)\textquotedblright{}.) Say it once more.
\end{tcolorbox}
\section[das Lesezeichen]{das Lesezeichen — bookmark}

\label{lesson:GE-C522-lesezeichen}



\begin{quote}
Before the new one: say the German for to borrow, then the German for to leaf.
\end{quote}

\subsection*{You'll want to know: das Lesezeichen}
\textbf{das Lesezeichen} — \textquotedblleft{}bookmark\textquotedblright{}.

\begin{quote}
Als Lesezeichen nimmt er eine alte Fahrkarte.
\end{quote}

\begin{quote}
\emph{He uses an old ticket as a bookmark.}
\end{quote}

Plural: \emph{die Lesezeichen}. Also a bookmark in a web browser.

\begin{cousinweb}
\emph{Lesezeichen} is built on \emph{lesen} \textquotedblleft{}to read\textquotedblright{} and \emph{Zeichen} \textquotedblleft{}sign\textquotedblright{}, inherited from Old High German \emph{zeihhan}, from Proto-Germanic \emph{taikną} (EtymDB).
\end{cousinweb}

\begin{grammarlens}[title={What you've built}]
One more word for telling, explaining and describing at B1.
\end{grammarlens}

\subsection*{Your turn}
\begin{itemize}
\raggedright
  \item \emph{Say it:} \emph{das Lesezeichen}
  \item \emph{Say it:} \emph{das Lesezeichen}, once more
  \item \emph{Say it:} a sentence of your own with \emph{das Lesezeichen}
  \item \emph{Recall:} say the German for to borrow, then the German for to leaf, then say \emph{das Lesezeichen} again
\end{itemize}

\begin{tcolorbox}[breakable,colback=teal!4,colframe=teal!35!black,title={Before you move on}]
What does das Lesezeichen mean? (\textquotedblleft{}Bookmark\textquotedblright{}.) Say it once more.
\end{tcolorbox}
\section[der Verlag]{der Verlag — publisher, publishing house}

\label{lesson:GE-C522-verlag}



\begin{quote}
Before the new one: say the German for to leaf, then the German for bookmark.
\end{quote}

\subsection*{You'll want to know: der Verlag}
\textbf{der Verlag} — \textquotedblleft{}publisher, publishing house\textquotedblright{}.

\begin{quote}
Ihr erstes Buch erscheint bei einem kleinen Verlag.
\end{quote}

\begin{quote}
\emph{Her first book is coming out with a small publisher.}
\end{quote}

Plural: \emph{die Verlage}. Built on \emph{verlegen} in the sense to publish.

\begin{cousinweb}
\emph{Verlag} is built on \emph{verlegen} \textquotedblleft{}to publish\textquotedblright{}, from \emph{ver-} and \emph{legen} \textquotedblleft{}to lay\textquotedblright{}, which is inherited from Old High German \emph{leggen}, from Proto-Germanic \emph{lagjaną} (EtymDB).
\end{cousinweb}

\begin{grammarlens}[title={What you've built}]
One more word for telling, explaining and describing at B1.
\end{grammarlens}

\subsection*{Your turn}
\begin{itemize}
\raggedright
  \item \emph{Say it:} \emph{der Verlag}
  \item \emph{Say it:} \emph{der Verlag}, once more
  \item \emph{Say it:} a sentence of your own with \emph{der Verlag}
  \item \emph{Recall:} say the German for to leaf, then the German for bookmark, then say \emph{der Verlag} again
\end{itemize}

\begin{tcolorbox}[breakable,colback=teal!4,colframe=teal!35!black,title={Before you move on}]
What does der Verlag mean? (\textquotedblleft{}Publisher, publishing house\textquotedblright{}.) Say it once more.
\end{tcolorbox}
\section[die Übersetzung]{die Übersetzung — translation}

\label{lesson:GE-C522-uebersetzung}



\begin{quote}
Before the new one: say the German for bookmark, then the German for publisher.
\end{quote}

\subsection*{You'll want to know: die Übersetzung}
\textbf{die Übersetzung} — \textquotedblleft{}translation\textquotedblright{}.

\begin{quote}
Ich lese das Buch in einer deutschen Übersetzung.
\end{quote}

\begin{quote}
\emph{I'm reading the book in a German translation.}
\end{quote}

Plural: \emph{die Übersetzungen}. \emph{Übersetzen} — to translate.

\begin{cousinweb}
\emph{Übersetzung} is built on \emph{übersetzen} \textquotedblleft{}to translate\textquotedblright{}, from \emph{über} \textquotedblleft{}across\textquotedblright{} and \emph{setzen} \textquotedblleft{}to set\textquotedblright{}, with \emph{-ung}; EtymDB lists it without a source.
\end{cousinweb}

\begin{grammarlens}[title={What you've built}]
One more word for telling, explaining and describing at B1.
\end{grammarlens}

\subsection*{Your turn}
\begin{itemize}
\raggedright
  \item \emph{Say it:} \emph{die Übersetzung}
  \item \emph{Say it:} \emph{die Übersetzung}, once more
  \item \emph{Say it:} a sentence of your own with \emph{die Übersetzung}
  \item \emph{Recall:} say the German for bookmark, then the German for publisher, then say \emph{die Übersetzung} again
\end{itemize}

\begin{tcolorbox}[breakable,colback=teal!4,colframe=teal!35!black,title={Before you move on}]
What does die Übersetzung mean? (\textquotedblleft{}Translation\textquotedblright{}.) Say it once more.
\end{tcolorbox}
